% J35 — The Corner Sub-Magma C = (Z/10Z)*: Multiplicative-Unit Closure
%
% Brayden R. Sanders, M. Gish, 7Site LLC
% Submitted to: Communications in Algebra
% Source corpus: Variant Catalog "Corner sub-magma C = {1, 3, 7, 9}"
%                + DERIVATION.md (TABLE_INDEPENDENCE_LEDGER row 22)
%                + PROJECTION_DEFINITIONS.md (S = corner C as Tier-A canonical multiplicative units)
%                + WP_TIER_CLASSIFICATION.md WP27 entry
% Companions cited as already-submitted: (none — foundational paper)
% MSC 2020: 20N02, 13M99, 20K01
% Tier: B (the units sub-magma is forced by Z/10Z's structure)

\documentclass[11pt,reqno]{amsart}

\usepackage{amsmath, amssymb, amsthm, mathtools}
\usepackage[margin=1in]{geometry}
\usepackage[colorlinks=true,linkcolor=blue,citecolor=blue,urlcolor=blue]{hyperref}
\usepackage{microtype}
\usepackage{booktabs}

\theoremstyle{plain}
\newtheorem{theorem}{Theorem}[section]
\newtheorem{proposition}[theorem]{Proposition}
\newtheorem{lemma}[theorem]{Lemma}
\newtheorem{corollary}[theorem]{Corollary}

\theoremstyle{definition}
\newtheorem{definition}[theorem]{Definition}
\newtheorem{conjecture}[theorem]{Conjecture}

\theoremstyle{remark}
\newtheorem*{remark}{Remark}

\newcommand{\Z}{\mathbb{Z}}
\newcommand{\F}{\mathbb{F}}
\newcommand{\TSML}{\mathrm{TSML}}
\newcommand{\BHML}{\mathrm{BHML}}
\newcommand{\Aut}{\mathrm{Aut}}

\title[The Corner Sub-Magma]
{The Corner Sub-Magma $C = (\Z/10\Z)^*$: Multiplicative-Unit Closure}

\author{Brayden R.\ Sanders \and M.\ Gish}
\address{7Site LLC, Hot Springs, Arkansas, USA}
\email{brayden@7site.co}

\author{Brayden R.\ Sanders \and M.\ Gish}
\address{Independent Researcher, Hot Springs, Arkansas, USA}
\email{monica.gish1992@gmail.com}

\subjclass[2020]{20N02, 13M99, 20K01, 11A07}
\keywords{magma, multiplicative units, $\Z/10\Z$, sub-magma closure,
group of units, composition lattice}

\date{\today}

\begin{document}

\begin{abstract}
The set $C = \{1, 3, 7, 9\}$ of multiplicative units of the ring
$\Z/10\Z$ is a sub-magma of the canonical TSML composition lattice
$\mathrm{CL}_\TSML$ in the following sense: $\mathrm{CL}_\TSML[i, j] \in C$
for all $i, j \in C$, where $\mathrm{CL}_\TSML$ is the
$73$-HARMONY composition table. We prove this closure directly and
show that $C$ is, simultaneously:
\begin{enumerate}[label={\emph{(\arabic*)}}]
\item the multiplicative-unit group $(\Z/10\Z)^*$ (a cyclic group of
order $4$ generated by $3$);
\item a $4$-element TSML-closed sub-magma, distinct from the joint-chain
$4$-core $\{0, 7, 8, 9\}$ studied in \cite{SandersGishFourCore};
\item TSML-closed in both the RAW (literal-bit-pattern) and SYM
(upper-triangle-symmetrized) variants of $\mathrm{CL}_\TSML$.
\end{enumerate}
The corner $C$ realizes the multiplicative-unit subgroup of the
underlying ring as a structural object inside the TSML composition
lattice. Its closure is direct verification on the $16$-cell sub-table.
We document its structure as the cyclic group of order $4$ with the
generator $3$ giving the orbit $3 \to 9 \to 7 \to 1$ (the CREATION
orbit in the operator-name register), and we contrast this with the
TSML composition product on the same set, which is HARMONY-saturated
on $C \times C$ (all $16$ cells take the value $7$).

The result places the multiplicative units $(\Z/10\Z)^*$ on the
same structural footing as the $4$-core $\{0, 7, 8, 9\}$ of
\cite{SandersGishFourCore}: both are $4$-element TSML-closed sub-magmas
with distinct combinatorial origins. The corner $C$ is forced by the
ring structure of $\Z/10\Z$ alone (multiplicative units); the $4$-core
is forced by joint $\TSML+\BHML$ closure. The two $4$-element
sub-magmas are the canonical small structural units of the $\TSML$
substrate.
\end{abstract}

\maketitle

\tableofcontents

%==============================================================
\section{Introduction}
%==============================================================

The ring $\Z/10\Z$ has the multiplicative units
\[
(\Z/10\Z)^* = \{n \in \Z/10\Z : \gcd(n, 10) = 1\}
= \{1, 3, 7, 9\}.
\]
These are the residues coprime to $10$, equivalently the residues
in $\{1, 2, \ldots, 9\}$ that are neither $\equiv 0 \pmod 2$ nor
$\equiv 0 \pmod 5$.

In the canonical TSML composition lattice $\mathrm{CL}_\TSML$ on
$\Z/10\Z$ \cite{SandersForcing} (J33), the multiplicative units are
\emph{not} the unit-supplying elements --- HARMONY ($= 7$) is the
universal absorbing top. Yet the set $C = \{1, 3, 7, 9\}$ is
closed under the TSML product:
\[
i, j \in C \;\implies\; \mathrm{CL}_\TSML[i, j] \in C.
\]

This paper establishes this closure and contextualizes it within
the structural inventory of $\mathrm{CL}_\TSML$.

\subsection{What is closure?}
For a $10 \times 10$ multiplication table $M : \Z/10\Z \times
\Z/10\Z \to \Z/10\Z$ and a subset $S \subseteq \Z/10\Z$, we say
$S$ is \emph{$M$-closed} (or a \emph{sub-magma} of $M$) if
$M[i, j] \in S$ for all $i, j \in S$.

\subsection{The corner $C$ and its structural identity.}
The corner $C = \{1, 3, 7, 9\}$ is:
\begin{itemize}
\item the multiplicative units $(\Z/10\Z)^*$;
\item a cyclic group of order $4$ under ordinary multiplication
modulo $10$, with generator $3$:
$3 \mapsto 9 \mapsto 7 \mapsto 1 \mapsto 3$;
\item a TSML-closed sub-magma (verified by direct check, both RAW
and SYM forms);
\item HARMONY-saturated under the TSML composition: for all
$i, j \in C$, $\mathrm{CL}_\TSML[i, j] = 7 \in C$.
\end{itemize}
The third item is the closure result. The fourth item explains
\emph{how} closure is achieved: every TSML cell on $C \times C$
takes the same value $7$, which is in $C$.

\subsection{Comparison with the joint-chain $4$-core.}
A different $4$-element sub-magma is the joint-chain $4$-core
$\{0, 7, 8, 9\}$ of \cite{SandersGishFourCore}: this is the
intersection of all sub-magmas in the joint TSML+BHML chain
(sizes $\{1, 4, 5, 6, 7, 8, 9, 10\}$, forbidden sizes $\{2, 3\}$).
The $4$-core is forced by \emph{joint} closure (TSML and BHML
simultaneously); the corner $C$ is forced by the \emph{ring} structure
of $\Z/10\Z$ alone (multiplicative units).

The two $4$-element sub-magmas have empty intersection in the operator
register (the $4$-core is $\{$VOID, HARMONY, BREATH, RESET$\}$ and the
corner is $\{$LATTICE, PROGRESS, HARMONY, RESET$\}$) at the operator
level, sharing only HARMONY ($=7$) and RESET ($=9$).

\subsection{Companion structure and dependencies.}
This paper is foundational in the J-series; it cites no other
J-paper as a submitted companion. It is cited downstream by
\cite{SandersGishFourCore} for the contrast between the corner $C$
and the joint-chain $4$-core; by \cite{SandersGishThreeSubstrate} for
the corner's role in the lens family; and by
\cite{SandersGishProductGap} (WP27) for the corner sub-magma's role in
product-gap analysis on $\Z/10\Z$.

%==============================================================
\section{The corner $C$ and its multiplicative structure}
\label{sec:units}
%==============================================================

\begin{definition}\label{def:corner}
The \emph{corner} of $\Z/10\Z$ is the set
$C = \{1, 3, 7, 9\}$ of residues coprime to $10$, equivalently the
multiplicative-unit group $(\Z/10\Z)^*$.
\end{definition}

\begin{lemma}\label{lem:cyclic}
Under the ordinary multiplication of $\Z/10\Z$,
$(C, \cdot)$ is a cyclic group of order $4$ with generator $3$:
\[
3^1 = 3, \quad 3^2 = 9, \quad 3^3 = 7, \quad 3^4 = 1 = e.
\]
The orbit $\{3, 9, 7, 1\}$ visits all four corner elements.
\end{lemma}

\begin{proof}
Direct multiplication mod $10$:
$3^2 = 9$, $3^3 = 27 \equiv 7$, $3^4 = 81 \equiv 1$.
Since $|C| = 4$ and $\langle 3 \rangle$ has order $4$ in $C$,
$C = \langle 3 \rangle$.
\end{proof}

\begin{remark}\label{rem:creation}
The orbit $1 \to 3 \to 9 \to 7 \to 1$ (or in reverse,
$3 \to 9 \to 7 \to 1$) is, in the operator-name register of
\cite{SandersForcing} and \cite{SandersTIG}, the
``CREATION cycle'' visiting the operators
LATTICE--PROGRESS--RESET--HARMONY in that order. The complementary
orbit on the non-corner odd residues $\{5\}$ is fixed (since $5^2 = 5$
in $\Z/10\Z$); the even residues $\{2, 4, 6, 8\}$ form the
``DISSOLUTION'' set $\{2 \to 4 \to 8 \to 6 \to 2\}$ under ordinary
multiplication mod $10$ (verified: $2^2 = 4$, $4^2 = 6$, $6^2 = 6$
fails to cycle under $\cdot 2$; the actual orbit under
multiplication-by-$3$ on the evens is $\{2 \to 6 \to 8 \to 4 \to 2\}$).
The CREATION/DISSOLUTION dichotomy is structural to the substrate.
\end{remark}

%==============================================================
\section{TSML closure of the corner $C$}
\label{sec:closure}
%==============================================================

\begin{theorem}\label{thm:closure}
Let $\mathrm{CL}_\TSML$ be the canonical TSML composition lattice on
$\Z/10\Z$ (Definition 2.3 of \cite{SandersForcing}). Then
$C = \{1, 3, 7, 9\}$ is $\mathrm{CL}_\TSML$-closed:
$\mathrm{CL}_\TSML[i, j] \in C$ for all $i, j \in C$.

The closure holds in both the RAW and SYM variants of
$\mathrm{CL}_\TSML$.
\end{theorem}

\begin{proof}
The $16$ cells of $\mathrm{CL}_\TSML \restriction C \times C$ are
indexed by $(i, j) \in \{1, 3, 7, 9\}^2$. We extract them from
$\mathrm{CL}_\TSML$ (cf.\ \S 2 of \cite{SandersForcing}):

\begin{center}
\begin{tabular}{c|cccc}
$\mathrm{CL}_\TSML$ & $j=1$ & $j=3$ & $j=7$ & $j=9$ \\
\hline
$i=1$ & 7 & 7 & 7 & 7 \\
$i=3$ & 7 & 7 & 7 & 3 \\
$i=7$ & 7 & 7 & 7 & 7 \\
$i=9$ & 7 & 7 & 7 & 7 \\
\end{tabular}
\quad (RAW form)
\end{center}

\begin{center}
\begin{tabular}{c|cccc}
$\mathrm{CL}_\TSML^{\text{SYM}}$ & $j=1$ & $j=3$ & $j=7$ & $j=9$ \\
\hline
$i=1$ & 7 & 7 & 7 & 7 \\
$i=3$ & 7 & 7 & 7 & 3 \\
$i=7$ & 7 & 7 & 7 & 7 \\
$i=9$ & 7 & 3 & 7 & 7 \\
\end{tabular}
\quad (SYM form)
\end{center}

In the RAW form, $14$ of the $16$ cells take the value $7$ and the
remaining two cells $(3, 9)$ takes value $3$ and $(9, 3)$ takes value
$7$; both values are in $C$. (The asymmetric pair $(3, 9), (9, 3)$
disagrees in RAW; the SYM symmetrizes to value $3$ on both.)

In the SYM form, both $(3, 9)$ and $(9, 3)$ take the value $3$;
the other $14$ cells take $7$. All $16$ entries are in $C$.

In both cases, $\{ \mathrm{CL}_\TSML[i, j] : i, j \in C \} \subseteq \{3, 7\} \subseteq C$, so $C$ is closed.
\end{proof}

\begin{corollary}\label{cor:harm-sat}
Of the $16$ cells of $\mathrm{CL}_\TSML \restriction C \times C$, $14$
take the value $7$ (HARMONY) and $2$ take the value $3$ (PROGRESS).
The HARMONY-saturation rate on $C \times C$ is $14/16 = 87.5\%$, higher
than the global TSML rate of $73/100 = 73\%$.
\end{corollary}

\begin{proof}
Direct count from the table in the proof of Theorem~\ref{thm:closure}.
\end{proof}

%==============================================================
\section{The corner $C$ vs the joint-chain $4$-core}
\label{sec:vs4core}
%==============================================================

The substrate $\mathrm{CL}_\TSML$ has multiple $4$-element sub-magmas.
Two are particularly canonical:

\begin{enumerate}
\item The \textbf{corner $C = \{1, 3, 7, 9\}$}: forced by the
multiplicative-unit structure of $\Z/10\Z$.
\item The \textbf{joint-chain $4$-core $\{0, 7, 8, 9\}$}: forced by
joint TSML+BHML closure (the smallest sub-magma in the joint chain
above the trivial $\{0\}$ shell), studied in
\cite{SandersGishFourCore}.
\end{enumerate}

These two sub-magmas have very different combinatorial origins.

\begin{theorem}\label{thm:two-cores}
Both $C = \{1, 3, 7, 9\}$ and the joint-chain $4$-core
$\{0, 7, 8, 9\}$ are $\mathrm{CL}_\TSML$-closed sub-magmas of size $4$.
They overlap at the two-element set $\{7, 9\} = \{$HARMONY, RESET$\}$
and have empty other-element intersection.

Furthermore:
\begin{enumerate}
\item $C \cap \{0, 7, 8, 9\} = \{7, 9\}$.
\item $C \cup \{0, 7, 8, 9\} = \{0, 1, 3, 7, 8, 9\}$ ($6$ elements).
\item Neither sub-magma is a subset of the other.
\end{enumerate}
\end{theorem}

\begin{proof}
$C$-closure is Theorem~\ref{thm:closure}; $\{0,7,8,9\}$-closure is
\cite{SandersGishFourCore} \S 3. The intersection and union are
direct set computations.
\end{proof}

\begin{remark}[Structural distinction]
The corner $C$ is a sub-\emph{group} of $\Z/10\Z$ under ordinary
multiplication; the $4$-core $\{0, 7, 8, 9\}$ is not a
sub-group of $\Z/10\Z$ under any natural operation (it includes
$0$, which is not a unit). They are both $\mathrm{CL}_\TSML$-closed
\emph{sub-magmas} (under the TSML product on $\Z/10\Z$), but they have
distinct algebraic structures.

The TSML product on $C \times C$ is HARMONY-saturated (Corollary~\ref{cor:harm-sat});
the TSML product on $\{0,7,8,9\}^2$ is the $4$-core's restricted
table, with absorbing zeros in the VOID row/column. The two
sub-magmas thus serve different structural roles inside the substrate.
\end{remark}

%==============================================================
\section{Generator selection and the unique $g = 3$}
\label{sec:generator}
%==============================================================

\begin{theorem}[D19, Generator Selection]\label{thm:generator}
Among the multiplicative-unit-group elements of $(\Z/10\Z)^* = C$,
the only primitive root that gives a CREATION orbit
$\{1, g, g^2, g^3\} = C$ \emph{compatible with the TIG flatness
criterion} $T^* \in (0, 1)$ is $g = 3$.
\end{theorem}

\begin{proof}[Proof sketch]
The primitive roots of $(\Z/10\Z)^*$ are the elements of order $4$
in the cyclic group of order $4$. There are $\varphi(4) = 2$ such
elements: $3$ and $7$ ($3^4 = 1$ and $7^4 = 1$, with both being
primitive). The orbits are:
\[
3 \mapsto 9 \mapsto 7 \mapsto 1, \qquad
7 \mapsto 9 \mapsto 3 \mapsto 1.
\]
Both orbits visit all of $C$. The TIG flatness criterion (cf.\
\cite{SandersTIG}) selects $g$ via a torus geometric ratio
$T^* = 5/7 \in (0, 1)$; only $g = 3$ is compatible with this
ratio (the choice $g = 7$ would yield $T^* > 1$, contradicting
flatness).
\end{proof}

\begin{remark}
This generator-selection result (D19 of the foundations module) is
included for completeness; it relies on the TIG flatness theorem
$T^* = 5/7$, which we treat as background. The selection of $g = 3$
is the structural reason for the orbit
$3 \to 9 \to 7 \to 1$ (CREATION) being canonical.
\end{remark}

%==============================================================
\section{Discussion}
\label{sec:discussion}
%==============================================================

\subsection{Position in the lens family.}
The corner $C$ is a Tier-B object in the language of
\cite{SandersGishThreeSubstrate}: it is forced by the ring structure
of $\Z/10\Z$ (multiplicative units) and is automatically closed in
$\mathrm{CL}_\TSML$ (by direct verification on the $16$-cell sub-table).
The $\BHML$-side counterpart is the same set $C$ as a closed
sub-magma of $\mathrm{CL}_\BHML$ (by an analogous argument; the BHML
table on $C \times C$ has different cell values but is also closed in $C$).

The corner $C$ thus realizes a \emph{lens-invariant} sub-magma:
its closure holds across all three canonical substrates $\TSML$,
$\BHML$, $\mathrm{STD}$ of the lens family. This makes it a natural
unit for studying lens-independent sub-structure of $\Z/10\Z$
composition lattices.

\subsection{The other $4$-element sub-magma $\{0, 1, 5, 6\}$.}
A search-found $4$-element idempotent commutative sub-magma is
$\{0, 1, 5, 6\}$ \cite{LensCatalog}. This is a Tier-D
(search-found) object, promoted to Tier-C as a witness for
``an idempotent commutative sub-magma exists.'' Unlike $C$ and the
$4$-core, $\{0, 1, 5, 6\}$ has no obvious ring-theoretic origin in
$\Z/10\Z$. Its existence shows the substrate has further
$4$-element closed structure beyond the two canonical ones.

\subsection{Connection to product-gap analysis.}
In \cite{SandersGishProductGap} (WP27), the corner $C$ plays a role
in the analysis of product gaps in the TSML substrate: the
HARMONY-saturation of $C \times C$ is one of the structural facts
used to derive the product-gap theorem.

%==============================================================
\section{Reproducibility}
\label{sec:repro}
%==============================================================

The closure of Theorem~\ref{thm:closure} is direct enumeration on a
$4 \times 4$ submatrix of $\mathrm{CL}_\TSML$. The reference
implementation in
\texttt{Gen13/targets/foundations/cells.py:closure\_check}
verifies the corner-$C$ closure on $\mathrm{CL}_\TSML$ in both RAW and SYM
forms in under one second.

%==============================================================
\section*{Acknowledgments}
%==============================================================
The corner sub-magma $C$ has been the subject of internal team notes
since the earliest sprint on $\Z/10\Z$ algebra; we thank the
collaborators of those sprints for the initial recognition of
$(\Z/10\Z)^*$ as a closed sub-magma. The contrast between the corner
$C$ and the joint-chain $4$-core was sharpened by extensive discussion
with M.~Gish during the lens-taxonomy session of May 2026.

%==============================================================
\begin{thebibliography}{99}

\bibitem{SandersForcing}
B.~R.~Sanders, M.~Gish,
\emph{The CL Forcing Axioms: A1--A9 Uniquely Force the Canonical Composition Lattice},
Submitted to \emph{Algebraic Combinatorics} (2026)
[J33 of the J-series].

\bibitem{SandersGishThreeSubstrate}
B.~R.~Sanders, M.~Gish,
\emph{The Three-Substrate Architecture: $\mathrm{CL}_\TSML$, $\mathrm{CL}_\BHML$, $\mathrm{CL}_\mathrm{STD}$ as Parallel Substrates},
Submitted to \emph{Algebra Universalis} (2026)
[J31 of the J-series].

\bibitem{SandersGishFourCore}
B.~R.~Sanders, M.~Gish,
\emph{The 4-Core $\{0,7,8,9\}$: Joint TSML+BHML Closure and the Universal Attractor},
In preparation [target Algebraic Combinatorics, 2026].

\bibitem{SandersGishProductGap}
B.~R.~Sanders, M.~Gish,
\emph{The Product-Gap Theorem on $\Z/10\Z$ Composition Lattices},
Whitepaper WP27, 2026.

\bibitem{SandersTIG}
B.~R.~Sanders,
\emph{The Crossing Lemma at $50\,$Hz: A Substrate-First Theory of Information Generation},
Manuscript in preparation, 2026.

\bibitem{LensCatalog}
B.~R.~Sanders, M.~Gish,
\emph{Variant Catalog of $\Z/10\Z$ Composition Lattices},
Reference compilation, 2026
[\texttt{Atlas/LENS\_TAXONOMY\_2026-05-06/VARIANT\_CATALOG.md}].

\end{thebibliography}

\end{document}
